\titleformat{\chapter}[display]
  {\normalfont\bfseries}{}{0pt}{\Large}
\chapter{Resumen}
En el presente trabajo se evaluó el potencial del fulereno C\textsubscript{60} funcionalizado con cuatro moléculas de glicina, C\textsubscript{60}(Gly)\textsubscript{4}, como posible inhibidor de la proteasa NS2B-NS3 del virus del dengue serotipo 2 (DENV-2). La elección de este sistema se basa en el interés por los derivados de fulereno con actividad antioxidante y en que su funcionalización con aminoácidos favorece su solubilidad en medios biológicos.

La caracterización electrónica, basada en la teoría del funcional de la densidad (DFT) y en las funciones de Fukui, permitió analizar la distribución de los grupos de glicina alrededor de la jaula del fulereno y los cambios asociados en su estructura electrónica. Las simulaciones de acoplamiento molecular mostraron afinidades favorables en regiones reportadas como sitios alostéricos y ortostéricos de la proteína. El análisis estructural identificó interacciones intermoleculares, entre ellas puentes de hidrógeno e interacciones de van der Waals. La presencia de los grupos de glicina modificó el potencial electrostático del aducto y su interacción con el entorno proteico.

La simulación de dinámica molecular del complejo proteína--aducto mostró estabilidad estructural durante el periodo analizado de 100 ns. Asimismo, la caracterización electrónica indicó una mayor facilidad para donar electrones respecto al fulereno sin funcionalizar. Estos resultados aportan información sobre el comportamiento del sistema en las condiciones simuladas y respaldan su estudio como candidato a inhibidor. Su actividad antiviral y su capacidad antioxidante requieren confirmación experimental.
